\documentclass[12pt]{article}

\title{Euclidean distance}
\author{
	Evan M Drake \\
	\texttt{drakeev@butte.edu}
}


\usepackage{snow-preamble}

%%%%%%%%%%%%%%%%%%%%%%%%%%%%%%%%%%%%
\begin{document}
	\maketitle

	%% body start
	\section{Rules}
		\indent In this assignment you will be creating a program to  measure Euclidean distance.
		You may know of this process (though not by this name) through use of the Pythagorean theorem.
		For this assignment you will need to devise a way to store a Cartesian coordinate;
			a function to perform the  Pythagorean theorem;
			and finally a way to link these actions together.
		You may not use any libraries for this assignment.
		There must be a write up turned in with this assignment explaining: your thinking,
			difficulties you had,
			and solutions to those challenges.

	\section{Difficulties}
		\indent For this assignment I'm trying something new.
		There will be different difficulties to this assignment that will gain you different points.
		\subsection{Difficulty one}
			\indent This is a simple Pythagorean theorem function.
			You only need make a function that executes the Pythagorean theorem on two Cartesian points.
			This will net you approx \verb|75%| score.

		\subsection{Difficulty two}
			\indent The full completion of this assignment on the integers.
			This will net you approx \verb|100%| score.

		\subsection{Difficulty three}
			\indent The full completion of this assignment on the real numbers.
			This will net you approx \verb|120%| score.

		\subsection{Difficulty four}
			\indent The full completion of this assignment on the complex numbers.
			This will net you approx \verb|150%| score.

	\section{Turn-in}
		\indent You will turn in a python file,
			for example \textit{euclid.py}.
		This file will contain the python logic you have constructed to solve the problem above (to some level).
		The second item of your turn in is an English write-up of your thoughts, challenges, and solutions during the assignment.
		Note that the English portion should be in the form of a \verb|.pdf| file.

	%% body finish

	\nocite{*}

	\newpage % bib
	\backmatter
	\addcontentsline{toc}{section}{References}
	\printbibliography

\end{document}

